\documentclass[10pt]{beamer}
\usetheme{default}
\useoutertheme{infolines}

\usepackage{amsmath}
\usepackage{amssymb}
\usepackage{mathtools}
\usepackage{graphicx}
\usepackage{listings}
\usepackage{hyperref}
\usepackage{tikz}
\usetikzlibrary{shapes,arrows,positioning,calc}

% Graphic search paths
\graphicspath{{../book/Figures/}{../../trans4mom/research/}{../../trans4mom/docs/Figures/}}

% Understated caption sizes
\setbeamerfont{caption}{size=\tiny}
\setbeamerfont{caption name}{size=\tiny}

% Academic Teal Palette
\definecolor{lightteal}{RGB}{236, 246, 246}
\definecolor{mediumteal}{RGB}{115, 186, 186}
\definecolor{darkteal}{RGB}{0, 110, 110}
\definecolor{highlightcolor}{RGB}{0, 80, 120}
\definecolor{alertred}{RGB}{160, 40, 40}

\setbeamercolor{background canvas}{bg=white}
\setbeamercolor{normal text}{fg=black}
\setbeamercolor{alerted text}{fg=alertred}
\setbeamercolor{frametitle}{bg=darkteal, fg=white}
\setbeamercolor{block title}{bg=darkteal, fg=white}
\setbeamercolor{block body}{bg=lightteal, fg=black}
\setbeamercolor{block title alerted}{bg=alertred, fg=white}
\setbeamercolor{block body alerted}{bg=lightteal, fg=black}
\setbeamercolor{block title example}{bg=mediumteal, fg=white}
\setbeamercolor{block body example}{bg=lightteal, fg=black}
\setbeamercolor{structure}{fg=darkteal}

% Python code listing styling
\lstset{
  language=Python,
  backgroundcolor=\color{lightteal},
  basicstyle=\ttfamily\tiny,
  keywordstyle=\color{darkteal}\bfseries,
  commentstyle=\color{gray}\itshape,
  stringstyle=\color{highlightcolor},
  breaklines=true,
  showstringspaces=false,
  frame=single,
  rulecolor=\color{mediumteal}
}

\title{The Foundations of Momentum and the Transition to AI}
\subtitle{A Three-Hour Masterclass on Systematic Trend Following, Deep Learning, and Transformers}
\author{Nick Firoozye}
\institute{UCL Computer Science \& Financial Computing}
\date{\today}

\begin{document}

% Title slide
\begin{frame}
  \titlepage
\end{frame}

% Masterclass Overview
\begin{frame}{Masterclass Syllabus}
  \begin{columns}[t]
    \begin{column}{0.48\textwidth}
      \textbf{Day 1: Foundations \& Shift to AI (1.5 h)}
      \begin{itemize}\footnotesize
        \item \textbf{Part 1 (45 mins): Anatomy of Momentum}
        \begin{itemize}\scriptsize
          \item Philosophy \& behavioural origins
          \item Classical TSMOM formulation
          \item Crisis alpha \& systemic shocks
          \item Limitations of fixed heuristics
        \end{itemize}
        \item \textbf{Part 2 (45 mins): Shift to Deep Learning}
        \begin{itemize}\scriptsize
          \item Deep Momentum Networks (DMNs)
          \item Direct optimization of Sharpe \& Sortino
          \item Turnover penalization \& differentiability
          \item Day 1 Review \& Discussion
        \end{itemize}
      \end{itemize}
    \end{column}
    \begin{column}{0.48\textwidth}
      \textbf{Day 2: Transformers \& Practice (1.5 h)}
      \begin{itemize}\footnotesize
        \item \textbf{Part 1 (45 mins): Momentum Transformer}
        \begin{itemize}\scriptsize
          \item Trans4mom architecture (VSN, LSTM, Attention)
          \item Inductive bias \& residual prior modulation
          \item Attention interpretability \& regime decoding
        \end{itemize}
        \item \textbf{Part 2 (45 mins): Trends, Data \& Practice}
        \begin{itemize}\scriptsize
          \item Modern architectures (TFT, PatchTST, Informer)
          \item Data reality: continuous futures rolls \& pitfalls
          \item Execution costs, sizing \& memory walls ($O(T^2)$)
          \item Day 2 Synthesis \& Q\&A
        \end{itemize}
      \end{itemize}
    \end{column}
  \end{columns}
\end{frame}

% =========================================================================
% DAY 1 - PART 1
% =========================================================================
\section{Day 1 - Part 1: Classical Momentum \& Trend Following}

\begin{frame}{The Philosophy of Momentum: Buy High, Sell Higher}
  \begin{block}{Deconstructing the Core Heuristic}
    Standard finance theory posits that asset returns follow a random walk with zero forecastable drift in frictionless markets. Momentum strategies invert the classical value prescription: rather than buying low and selling high, trend following executes the rule \emph{buy high and sell higher, or sell short and cover lower}.
  \end{block}
  \vspace{0.2cm}
  \begin{columns}[t]
    \begin{column}{0.48\textwidth}
      \textbf{Behavioural Underpinnings}
      \begin{itemize}\scriptsize
        \item \textbf{Disposition effect:} Investors realise gains prematurely while holding losers, causing initial underreaction to fundamental news.
        \item \textbf{Anchoring \& conservatism:} Slow Bayesian updating across market participants creates auto-correlated price discovery.
        \item \textbf{Institutional herding:} Benchmark tracking and delegated portfolio mandates produce persistent capital flows.
      \end{itemize}
    \end{column}
    \begin{column}{0.48\textwidth}
      \textbf{Microstructural Mechanics}
      \begin{itemize}\scriptsize
        \item \textbf{Liquidity provision limits:} Capital constraints among market makers prevent instantaneous price convergence.
        \item \textbf{Feedback loops:} Systematic rebalancing and stop-loss execution amplify directional displacement over intermediate horizons.
        \item \textbf{Multi-horizon diffusion:} Information diffuses at differing speeds across diverse participant mandates.
      \end{itemize}
    \end{column}
  \end{columns}
\end{frame}

\begin{frame}{Two Flavours of Momentum: TSMOM vs. XSMOM}
  \begin{block}{Time-Series Momentum (TSMOM / Absolute Momentum)}
    Evaluates an asset strictly against its own historical trajectory:
    \begin{equation}
      s_{i,t}^{\text{TS}} = \text{sign}\left(r_{i,t}^{(k)}\right), \quad r_{i,t}^{(k)} = \ln\left(\frac{P_{i,t}}{P_{i,t-k}}\right)
    \end{equation}
    Dominant in managed futures and CTAs across liquid multi-asset derivatives (bonds, rates, FX, equity index futures, commodities). Documented comprehensively by \href{https://doi.org/10.1016/j.jfineco.2011.11.003}{Moskowitz, Ooi, and Pedersen (2012)}.
  \end{block}
  \vspace{0.1cm}
  \begin{block}{Cross-Sectional Momentum (XSMOM / Relative Strength)}
    Evaluates an asset relative to the contemporaneous cross-sectional distribution of peer returns:
    \begin{equation}
      s_{i,t}^{\text{XS}} = \mathbf{1}_{\{r_{i,t}^{(k)} \in Q_{\text{top}}\}} - \mathbf{1}_{\{r_{i,t}^{(k)} \in Q_{\text{bottom}}\}}
    \end{equation}
    Dominant in quantitative equity factor investing. Documented originally by \href{https://doi.org/10.1111/j.1540-6261.1993.tb04702.x}{Jegadeesh and Titman (1993)} and extended by \href{https://doi.org/10.1111/jofi.12021}{Asness, Moskowitz, and Pedersen (2013)}.
  \end{block}
\end{frame}

\begin{frame}{Classical Signal Formulations \& The Filter Equivalence Theorem}
  \begin{block}{Digital Trend Filtering Formulations}
    \begin{itemize}\scriptsize
      \item \textbf{SMA Crossover:} $s_t = \text{sign}\left(\text{SMA}_t(k_{\text{fast}}) - \text{SMA}_t(k_{\text{slow}})\right)$
      \item \textbf{EWMA Crossover:} $s_t = \text{sign}\left(\text{EWMA}_t(\lambda_{\text{fast}}) - \text{EWMA}_t(\lambda_{\text{slow}})\right)$
      \item \textbf{Donchian Channel:} $s_t = +1 \text{ if } P_t = \max_{j \in [0, k]} P_{t-j}, \quad -1 \text{ if } P_t = \min_{j \in [0, k]} P_{t-j}$
    \end{itemize}
  \end{block}
  \vspace{0.1cm}
  \begin{block}{The Filter Equivalence Theorem (\href{https://ssrn.com/abstract=2465623}{Bruder \& Gaussel 2011}, \href{http://www.thierry-roncalli.com/download/Alternative_Risk_Premia.pdf}{Hamdan et al. 2012})}
    Any linear trend signal can be expressed as a linear filter of historical returns with lag weights $\phi_k$:
    \begin{equation}
      s_{t-1} = \sum_{k=1}^\infty \phi_k r_{t-k}
    \end{equation}
    Different digital filters (SMA, EWMA, MACD) yield strategy returns with mutual correlations exceeding $0.90$. The apparent diversity of classical trend heuristics collapses mathematically to an implicit weighting over autocorrelation lag coefficients.
  \end{block}
\end{frame}

\begin{frame}{Multi-Century Empirical Persistence: The 200-Year Evidence}
  \begin{block}{Long-Horizon Statistical Proof (\href{https://arxiv.org/abs/1404.7797}{Lempérière et al. 2014})}
    Examining two centuries of multi-asset data (1800 to 2014) confirms that time-series momentum has delivered a persistent annualized Sharpe ratio of $\approx 0.70$ through agricultural, industrial, and digital economic eras.
  \end{block}
  \begin{figure}[ht]
    \centering
    \includegraphics[width=0.55\linewidth]{L5/Trend_Following_200y.png}
    \caption{Two centuries of trend following across equity, bond, FX, and commodity futures (Lempérière et al. 2014).}
  \end{figure}
\end{frame}

\begin{frame}{The Quant Meltdown Survival: Crisis Alpha}
  \begin{block}{Portfolio Stabilisation During Systemic Shocks}
    Institutional allocators incorporate trend-following strategies because they generate convex positive returns precisely when long-only equity and credit portfolios suffer severe drawdowns.
  \end{block}
  \begin{figure}[ht]
    \centering
    \includegraphics[width=0.65\linewidth]{L1/sp500_vs_cta_annual.png}
    \caption{Annual returns during major crisis regimes: S\&P 500 equity index versus CTA benchmark.}
  \end{figure}
\end{frame}

\begin{frame}{Mathematical Convexity \& Lookback Straddle Equivalence}
  \begin{block}{Fung-Hsieh Lookback Option Equivalence (\href{https://people.duke.edu/~dah7/RFS2001.pdf}{Fung \& Hsieh 2001})}
    A continuous-time trend-following strategy with dynamic stop-losses and position flipping generates an asymmetric return distribution mathematically equivalent to a portfolio of lookback straddles:
    \begin{equation}
      \text{Payoff}_{\text{lookback}} = \max_{0 \le t \le T} P_t - \min_{0 \le t \le T} P_t
    \end{equation}
  \end{block}
  \vspace{0.1cm}
  \begin{block}{The Martin-Zou Positive Skewness Theorem (\href{https://www.risk.net/derivatives/structured-products/2194247/momentum-trading-skews-me}{Martin \& Zou 2012})}
    When filtering i.i.d. Gaussian innovations $X_t \sim \mathcal{N}(0, \sigma^2)$ through an exponentially weighted moving average filter $s_t = c \sum_{k=1}^\infty \rho^k X_{t-k}$, the compounded strategy return possesses strictly positive skewness over horizon $T$:
    \begin{equation}
      \mathbb{E}\left[R_{t+T}^3\right] = 6 \sum_{j=1}^T (T - j) a_j \gamma_j^a > 0
    \end{equation}
    Unlike carry and short-volatility strategies that suffer severe negative skewness (tail crashes), momentum delivers structural positive skewness.
  \end{block}
\end{frame}

\begin{frame}{Empirical Evidence of Asymmetric Skewness}
  \begin{columns}[c]
    \begin{column}{0.48\textwidth}
      \begin{figure}[ht]
        \centering
        \includegraphics[width=0.95\linewidth]{L5/chf_usd_skew.png}
        \caption{Right-skewed return distribution of trend following (Potters \& Bouchaud).}
      \end{figure}
    \end{column}
    \begin{column}{0.48\textwidth}
      \begin{figure}[ht]
        \centering
        \includegraphics[width=0.95\linewidth]{L5/MartinsSkewnessEWMA1.png}
        \caption{Analytical positive skewness vs. investment horizon (Martin \& Zou 2012).}
      \end{figure}
    \end{column}
  \end{columns}
\end{frame}

\begin{frame}{Limitations of Classical Heuristics: Momentum Crashes}
  \begin{block}{The Rebound Crash Anomaly (\href{https://www.nber.org/papers/w20439}{Daniel \& Moskowitz 2016})}
    Unhedged momentum strategies suffer devastating crashes during sudden market inflection points, particularly V-shaped recoveries following panics (e.g. July--August 1932, March--April 2009). The short leg behaves like a written call option with extreme negative gamma on distressed recovery assets.
  \end{block}
  \begin{figure}[ht]
    \centering
    \includegraphics[width=0.55\linewidth]{L5/XSMomDDGFC.png}
    \caption{Cross-sectional momentum drawdown during the 2008--2009 GFC rebound (Daniel \& Moskowitz 2016).}
  \end{figure}
\end{frame}

\begin{frame}{Why Fixed Heuristics Break Down}
  \begin{columns}[t]
    \begin{column}{0.48\textwidth}
      \textbf{1. The Lag Drag Dilemma}
      \begin{itemize}\scriptsize
        \item Long lookback windows (e.g., 200 days) protect against high-frequency noise but bleed capital when regimes pivot.
        \item Short windows (e.g., 10 days) adapt rapidly to pivots but suffer severe transaction cost bleed in choppy, sideways markets.
      \end{itemize}
      \vspace{0.1cm}
      \textbf{2. Fixed-Threshold Brittleness}
      \begin{itemize}\scriptsize
        \item Discrete binary decisions ($s_t \in \{-1, +1\}$) ignore continuous uncertainty and regime transitions.
        \item Linear scaling fails to model cross-asset volatility spillovers and correlation shocks.
      \end{itemize}
    \end{column}
    \begin{column}{0.48\textwidth}
      \textbf{3. Classical State-Space Extensions}
      \begin{itemize}\scriptsize
        \item \textbf{Kalman filtering:} Estimates latent trend slope dynamically, but assumes Gaussian noise innovations:
        \begin{align*}
          y_t &= L_t + \epsilon_t, \quad \epsilon_t \sim \mathcal{N}(0, \sigma_\epsilon^2) \\
          L_t &= L_{t-1} + T_{t-1} + \eta_{1,t} \\
          T_t &= T_{t-1} + \eta_{2,t}
        \end{align*}
        \item \textbf{Bayesian Changepoint Detection (BOCPD):} Computes run-length hazard rates, but struggles with multi-asset feature interactions.
      \end{itemize}
      \vspace{0.1cm}
      \textbf{Conclusion:} Setting the stage for non-linear, deep learning architectures.
    \end{column}
  \end{columns}
\end{frame}

% =========================================================================
% DAY 1 - PART 2
% =========================================================================
\section{Day 1 - Part 2: Transitioning to Deep Learning}

\begin{frame}{From Fixed Rules to Deep Momentum Networks (DMNs)}
  \begin{block}{The Core Machine Learning Paradigm Shift (\href{https://arxiv.org/abs/1904.04912}{Lim, Zohren, \& Roberts 2019})}
    Traditional systematic pipelines decompose the trading problem into distinct serial modules: signal generation, volatility estimation, and portfolio optimization. Deep Momentum Networks combine these disjoint steps into an end-to-end differentiable computational graph.
  \end{block}
  \vspace{0.2cm}
  \begin{columns}[t]
    \begin{column}{0.48\textwidth}
      \textbf{Classical Pipeline}
      \begin{itemize}\scriptsize
        \item Heuristic price filtering (e.g. moving average difference).
        \item Static sign or linear scaling.
        \item Disjoint ex-post volatility weighting ($w_t = \sigma_{\text{target}} / \hat{\sigma}_t$).
        \item Unaware of downstream slippage and turnover costs during signal construction.
      \end{itemize}
    \end{column}
    \begin{column}{0.48\textwidth}
      \textbf{Deep Momentum Network (DMN)}
      \begin{itemize}\scriptsize
        \item Inputs: multi-horizon returns and normalized volatility features across asset universes.
        \item Architecture: Recurrent layers (LSTM) learn path-dependent memory states.
        \item Direct output: continuous position allocations $z_{i,t} \in (-1, 1)$ bounded by activation functions (e.g. $\tanh$).
      \end{itemize}
    \end{column}
  \end{columns}
\end{frame}

\begin{frame}{Optimizing for the Right Metric: Why MSE Fails}
  \begin{block}{The Breakdown of Mean Squared Error in Trading}
    Supervised learning architectures trained on point forecasting (minimizing Mean Squared Error $\mathcal{L}_{\text{MSE}} = \frac{1}{N}\sum (\hat{r}_{t+1} - r_{t+1})^2$) exhibit near-zero performance when deployed in live strategy execution.
  \end{block}
  \vspace{0.2cm}
  \begin{columns}[t]
    \begin{column}{0.48\textwidth}
      \textbf{Why MSE Disconnects from PnL}
      \begin{itemize}\scriptsize
        \item \textbf{Low Signal-to-Noise Ratio:} In financial returns, variance is dominated by stochastic innovations. A model minimizing MSE converges to predicting zero return everywhere.
        \item \textbf{Symmetric Penalty:} MSE treats a missing upside rally identically to an inverse position during a crash, ignoring portfolio risk geometry.
        \item \textbf{Directional Indifference:} A small positive prediction when returns are large and positive gives low MSE, but fails to size capital aggressively.
      \end{itemize}
    \end{column}
    \begin{column}{0.48\textwidth}
      \textbf{The Financial Objective}
      \begin{itemize}\scriptsize
        \item Quant practitioners do not trade point forecasts; they trade risk-adjusted distributions.
        \item The model must solve for the optimal policy mapping history to position: $f_\theta: \mathcal{X}_{1:t} \to z_t$.
        \item The loss function must directly reflect the utility of the risk-adjusted portfolio return:
        \begin{equation*}
          \max_\theta \quad \text{Utility}(R_1, \dots, R_T)
        \end{equation*}
      \end{itemize}
    \end{column}
  \end{columns}
\end{frame}

\begin{frame}{Formulating Differentiable Sharpe and Sortino Loss Functions}
  \begin{block}{End-to-End Sharpe Loss}
    Let $z_t \in [-1, 1]$ denote the portfolio position at time $t$, and $r_{t+1}$ the realized volatility-scaled return. The strategy return net of transaction cost $c$ is:
    \begin{equation}
      R_{t+1} = z_t r_{t+1} - c \cdot |z_t - z_{t-1}|
    \end{equation}
    Over training batch sequence of length $B$, we compute empirical mean and variance:
    \begin{equation}
      \hat{\mu} = \frac{1}{B} \sum_{t=1}^B R_t, \quad \hat{\sigma}^2 = \frac{1}{B-1} \sum_{t=1}^B (R_t - \hat{\mu})^2 + \epsilon
    \end{equation}
    The loss function minimized during gradient backpropagation is the negative annualized Sharpe:
    \begin{equation}
      \mathcal{L}_{\text{Sharpe}}(\theta) = -\sqrt{252} \cdot \frac{\hat{\mu}}{\hat{\sigma}}
    \end{equation}
  \end{block}
\end{frame}

\begin{frame}{Differentiable Sharpe vs. Sortino Loss Formulations}
  \begin{block}{Penalizing Downside Semivariance: Differentiable Sortino Loss}
    To maintain positive skewness without penalizing upside windfalls, we can replace total standard deviation with downside semi-deviation:
    \begin{equation}
      \hat{\sigma}_{\text{down}}^2 = \frac{1}{B} \sum_{t=1}^B \min(R_t - \tau, 0)^2 + \epsilon
    \end{equation}
    where $\tau$ is the minimum acceptable return (MAR, typically 0). The Sortino loss is:
    \begin{equation}
      \mathcal{L}_{\text{Sortino}}(\theta) = -\sqrt{252} \cdot \frac{\hat{\mu} - \tau}{\hat{\sigma}_{\text{down}}}
    \end{equation}
  \end{block}
  \vspace{0.1cm}
  \begin{block}{Subgradient Smoothing for Turnover Cost}
    Because $|z_t - z_{t-1}|$ has a non-differentiable cusp at zero, production implementations apply smooth approximations:
    \begin{equation}
      ||\Delta z_t||_{\epsilon} \approx \sqrt{(z_t - z_{t-1})^2 + \epsilon_{\text{smooth}}} \quad \text{or} \quad \text{Huber}(\Delta z_t)
    \end{equation}
    enabling stable backward passes through AdamW optimizers.
  \end{block}
\end{frame}

\begin{frame}[fragile]{PyTorch Implementation: Differentiable Sharpe Loss}
\begin{lstlisting}
import torch
import torch.nn as nn

class DifferentiableSharpeLoss(nn.Module):
    def __init__(self, trans_cost: float = 0.0001, ann_factor: float = 252.0):
        super().__init__()
        self.trans_cost = trans_cost
        self.ann_factor = torch.sqrt(torch.tensor(ann_factor))

    def forward(self, positions: torch.Tensor, returns: torch.Tensor) -> torch.Tensor:
        # Align: positions at t generate return at t+1
        active_pos = positions[:, :-1]
        active_rets = returns[:, 1:]
        
        # Compute smooth transaction cost on turnover
        pos_delta = positions[:, 1:] - positions[:, :-1]
        costs = self.trans_cost * torch.sqrt(pos_delta.pow(2) + 1e-8)
        
        # Net portfolio return series
        net_returns = (active_pos * active_rets) - costs
        
        # Batch mean and standard deviation
        mean_ret = torch.mean(net_returns)
        std_ret = torch.std(net_returns, unbiased=True) + 1e-6
        
        # Return negative annualized Sharpe ratio
        sharpe = (mean_ret / std_ret) * self.ann_factor.to(returns.device)
        return -sharpe
\end{lstlisting}
\end{frame}

\begin{frame}{Day 1 Summary \& Discussion}
  \begin{block}{Key Takeaways from Day 1}
    \begin{itemize}\scriptsize
      \item \textbf{Economic Anchor:} Momentum is backed by two centuries of empirical persistence and exhibits structural positive return skewness (Martin-Zou theorem, lookback straddle equivalence).
      \item \textbf{Filter Equivalence:} All linear digital moving average filters produce highly correlated positions ($>0.90$).
      \item \textbf{Vulnerability:} Fixed heuristics suffer severe lag drag and momentum crashes during sharp V-bottom market inflection points (Daniel \& Moskowitz).
      \item \textbf{Objective Alignment:} Machine learning for trend following requires replacing MSE with end-to-end differentiable risk-adjusted metrics (Sharpe/Sortino loss) including explicit turnover transaction cost penalties.
    \end{itemize}
  \end{block}
  \vspace{0.2cm}
  \begin{alertblock}{Looking Forward to Day 2}
    While recurrent neural networks (LSTMs) capture local path dependencies, they struggle with multi-year regime shifts and non-stationarity. Tomorrow, we examine the \textbf{Momentum Transformer}, its attention mechanism, and practical deployment realities.
  \end{alertblock}
\end{frame}

% =========================================================================
% DAY 2 - PART 1
% =========================================================================
\section{Day 2 - Part 1: Inside the Momentum Transformer}

\begin{frame}{Why Transformers for Momentum? Addressing Non-Stationarity}
  \begin{block}{The Non-Stationarity Trilemma in Financial Time Series}
    Financial markets undergo structural regime changes driven by macroeconomic shocks, central bank liquidity shifts, and technological disruptions.
  \end{block}
  \vspace{0.2cm}
  \begin{columns}[t]
    \begin{column}{0.48\textwidth}
      \textbf{Limits of Classical Baselines}
      \begin{itemize}\scriptsize
        \item \textbf{Linear filters} assume future statistical moments mirror the past.
        \item \textbf{Recurrent networks (LSTMs)} store past dynamics in hidden states, but suffer exponential gradient decay over multi-quarter horizons.
        \item In an LSTM, historical regimes (e.g. 2008 GFC, 2020 pandemic) are progressively overwritten by recent observations.
      \end{itemize}
    \end{column}
    \begin{column}{0.48\textwidth}
      \textbf{The Transformer Advantage}
      \begin{itemize}\scriptsize
        \item \textbf{Associative memory:} Multi-head self-attention treats historical observations as an open, content-addressable database.
        \item When today's volatility shock or yield curve shift resembles an event from three years ago, attention heads reference that regime directly without path degradation.
      \end{itemize}
    \end{column}
  \end{columns}
\end{frame}

\begin{frame}{The Trans4mom Architecture (\href{https://arxiv.org/abs/2112.08534}{Wood, Giegerich, Roberts, \& Zohren 2022})}
  \begin{block}{Architectural Pipeline Overview}
    The Momentum Transformer decomposes the regime-inference process into four coordinated stages:
  \end{block}
  \vspace{0.2cm}
  \begin{columns}[t]
    \begin{column}{0.48\textwidth}
      \textbf{1. Variable Selection Network (VSN)}
      \begin{itemize}\scriptsize
        \item Sparse feature gating across multi-horizon inputs (returns, MACD scales, volatility, changepoint scores).
        \item Dynamic relevance weights conditioned on static asset class context.
      \end{itemize}
      \vspace{0.1cm}
      \textbf{2. Locality Enhancement (LSTM Layer)}
      \begin{itemize}\scriptsize
        \item Processes high-frequency sequential order and local volatility clustering.
      \end{itemize}
    \end{column}
    \begin{column}{0.48\textwidth}
      \textbf{3. Multi-Head Causal Self-Attention}
      \begin{itemize}\scriptsize
        \item Relates historical timesteps across horizons of 252 trading days.
        \item Enforces strict causal upper-triangular masking to prohibit future information leakage.
      \end{itemize}
      \vspace{0.1cm}
      \textbf{4. Residual Signal Head ($\alpha$-Modulation)}
      \begin{itemize}\scriptsize
        \item Modulates a classical baseline momentum prior rather than estimating unconstrained positioning.
      \end{itemize}
    \end{column}
  \end{columns}
\end{frame}

\begin{frame}{Variable Selection Networks (VSN) \& Static Context Conditioning}
  \begin{block}{Feature Sparsity \& Noise Suppression}
    Financial inputs suffer from severe noise. The VSN uses Gated Residual Networks (GRNs) to generate non-linear feature coefficients $v_t^{(j)}$:
    \begin{equation}
      v_t = \text{Softmax}\left(\text{GRN}_v(x_t, c_{\text{static}})\right)
    \end{equation}
    where $c_{\text{static}}$ is the asset class embedding (Equities, Fixed Income, Energy, Metals, FX). The processed representation combines selected features:
    \begin{equation}
      \tilde{x}_t = \sum_{j=1}^M v_t^{(j)} \cdot \text{GRN}_j\left(x_t^{(j)}\right)
    \end{equation}
  \end{block}
  \vspace{0.1cm}
  \begin{itemize}\scriptsize
    \item In calm bull markets, the VSN allocates weight to slow momentum features (e.g., 252-day return, MACD 32/96).
    \item During sudden market shocks, the VSN shifts focus toward fast return windows and Bayesian changepoint hazard metrics.
  \end{itemize}
\end{frame}

\begin{frame}{Structural Inductive Bias: The Residual Prior Anchoring}
  \begin{alertblock}{The Overfitting Hazard of Unconstrained Transformers}
    Unconstrained transformers possess enormous parameter capacity ($O(10^5 \text{ to } 10^7)$ weights). In low signal-to-noise financial environments, unguided transformers readily memorize spurious historical patterns, yielding catastrophic out-of-sample collapse.
  \end{alertblock}
  \vspace{0.2cm}
  \begin{block}{The Residual "Handcuffs" Principle}
    Rather than requiring the network to learn basic directional momentum from scratch, the Momentum Transformer anchors its output to a robust classical baseline prior (e.g. equal-weighted multi-scale MACD):
    \begin{equation}
      \text{Position}_t = \tanh\left(\text{Prior}_t \cdot (1 + \alpha_t)\right)
    \end{equation}
    where $\alpha_t \in \mathbb{R}$ is the learned residual modulation output.
  \end{block}
  \vspace{0.1cm}
  \begin{itemize}\scriptsize
    \item Zero-weight initialization of the final linear projection layer guarantees that $\alpha_t \to 0$ initially: the model starts execution as a proven classical CTA before fine-tuning regime adjustments.
    \item The network learns \emph{when to scale up, size down, or neutralize positions}, avoiding unanchored direction flipping.
  \end{itemize}
\end{frame}

\begin{frame}{Interpretability: Demystifying the Black Box}
  \begin{block}{Decoding Multi-Head Attention Weights}
    Unlike black-box multilayer perceptrons, transformer attention matrices provide an exact, inspectable record of historical temporal dependencies:
    \begin{equation}
      A_{h}(t, \tau) = \text{Softmax}\left(\frac{Q_h(t) K_h(\tau)^\top}{\sqrt{d_k}} + M_{\text{causal}}\right)
    \end{equation}
  \end{block}
  \vspace{0.1cm}
  \begin{columns}[c]
    \begin{column}{0.48\textwidth}
      \begin{figure}[ht]
        \centering
        \includegraphics[width=0.98\linewidth]{attention_ES_F.png}
        \caption{Empirical S\&P 500 futures attention weights across time lags.}
      \end{figure}
    \end{column}
    \begin{column}{0.48\textwidth}
      \textbf{Emergent Head Specialization}
      \begin{itemize}\scriptsize
        \item \textbf{Slow Trend Head:} Concentrates attention across 180--252 day lookback horizons, confirming persistent macro drift.
        \item \textbf{Mean-Reversion Head:} Places intense weight on recent 1--5 day shocks, scaling down positions after excessive short-term extensions.
        \item \textbf{Regime-Matching Head:} Attends to historical inflection points during market pullbacks.
      \end{itemize}
    \end{column}
  \end{columns}
\end{frame}

\begin{frame}{Cross-Sectional Attention Maps \& Strategy Blending}
  \begin{columns}[c]
    \begin{column}{0.5\textwidth}
      \begin{figure}[ht]
        \centering
        \includegraphics[width=0.98\linewidth]{cs_attention_map.png}
        \caption{Cross-asset attention map illustrating cross-sectional regime identification.}
      \end{figure}
    \end{column}
    \begin{column}{0.48\textwidth}
      \begin{block}{Interpreting Regime Signatures}
        When systemic market stress strikes:
        \begin{itemize}\scriptsize
          \item Attention across asset classes aligns synchronously onto shared systemic shock events.
          \item The model rapidly shifts weight away from slow moving averages toward fast changepoint indicators.
          \item Classical rule-based CTAs suffer lag drag during these pivots, whereas attention-driven modulation de-leverages early.
        \end{itemize}
      \end{block}
    \end{column}
  \end{columns}
\end{frame}

% =========================================================================
% DAY 2 - PART 2
% =========================================================================
\section{Day 2 - Part 2: New AI Trends, Data Bottlenecks, \& Implementation}

\begin{frame}{Alternative Transformer Architectures for Financial Time Series}
  \begin{block}{Comparative Survey of Modern Sequence Models}
    The time-series deep learning landscape has evolved rapidly beyond generic multi-head attention:
  \end{block}
  \vspace{0.2cm}
  \begin{columns}[t]
    \begin{column}{0.48\textwidth}
      \textbf{1. Temporal Fusion Transformer (TFT)}
      \begin{itemize}\scriptsize
        \item Developed by Lim et al. (2021).
        \item Separates static metadata, observed historical inputs, and known future inputs (e.g. calendar/rollover dates).
        \item Multi-quantile outputs for full return distribution modeling.
      \end{itemize}
      \vspace{0.1cm}
      \textbf{2. PatchTST (Nie et al. 2023)}
      \begin{itemize}\scriptsize
        \item Segments time series into contiguous sub-series "patches" (tokens).
        \item Drastically reduces sequence length from $L$ to $L/P$, enabling lookbacks of thousands of bars.
        \item Channel-independent processing preserves individual asset characteristics.
      \end{itemize}
    \end{column}
    \begin{column}{0.48\textwidth}
      \textbf{3. Informer \& Autoformer}
      \begin{itemize}\scriptsize
        \item \textbf{Informer (Zhou et al.):} ProbSparse attention selects top-$k$ dominant queries, reducing computation to $O(L \log L)$.
        \item \textbf{Autoformer (Wu et al.):} Replaces self-attention with auto-correlation mechanisms based on Wiener-Khinchin theorem.
      \end{itemize}
      \vspace{0.1cm}
      \begin{alertblock}{Practitioner Verdict}
        In finance, long horizons without inductive bias increase overfitting risk. Architectural parsimony combined with structural priors consistently outperforms massive unconstrained models.
      \end{alertblock}
    \end{column}
  \end{columns}
\end{frame}

\begin{frame}{The Data Reality Check: Why Model Architecture is Half the Battle}
  \begin{alertblock}{The Hard Truth of Financial Deep Learning}
    No neural network, regardless of parameter count or attention sophistication, can overcome corrupted, contaminated, or improperly adjusted data inputs.
  \end{alertblock}
  \vspace{0.2cm}
  \begin{columns}[t]
    \begin{column}{0.48\textwidth}
      \textbf{1. The Roll Gap Catastrophe}
      \begin{itemize}\scriptsize
        \item Futures contracts expire periodically (quarterly, monthly).
        \item Simplistic free data (e.g. unadjusted Yahoo continuous contracts) concatenates adjacent contracts without backward adjustment.
        \item Contango/backwardation creates artificial 5--10\% overnight jumps at expiry.
        \item These artificial price gaps distort EWMA volatility and inject spurious gradient updates into Sharpe loss backpropagation.
      \end{itemize}
    \end{column}
    \begin{column}{0.48\textwidth}
      \textbf{2. Roll Methodology Matters}
      \begin{itemize}\scriptsize
        \item \textbf{Panama (Difference) Adjustment:} Preserves absolute price differences; causes negative historical prices in commodities (e.g. Nat Gas, Crude).
        \item \textbf{Ratio (Percentage) Adjustment:} Preserves true compounding percentage returns across rolls. \emph{Mandatory for quantitative Sharpe models.}
        \item \textbf{Institutional Data Feeds:} Sourcing from professional vendors (Pinnacle Data, CSI, Databento, Bloomberg) or using total-return ETF proxies (SPY, TLT, GLD).
      \end{itemize}
    \end{column}
  \end{columns}
\end{frame}

\begin{frame}{The Breadth-Regularisation Link: 24 vs. 46 Assets}
  \begin{block}{Cross-Asset Breadth as Regularisation}
    High-capacity transformer models require adequate cross-sectional variance to learn generalized regime representations rather than memorizing individual asset histories.
  \end{block}
  \vspace{0.2cm}
  \begin{columns}[c]
    \begin{column}{0.48\textwidth}
      \begin{block}{Narrow Universe (24 Assets)}
        \begin{itemize}\scriptsize
          \item In-Sample Sharpe: $1.27$
          \item Out-of-Sample Sharpe: \textbf{-1.00}
          \item \emph{Diagnosis:} Severe overfitting. The model memorized asset-specific noise patterns.
        \end{itemize}
      \end{block}
    \end{column}
    \begin{column}{0.48\textwidth}
      \begin{block}{Broad Universe (46 Assets)}
        \begin{itemize}\scriptsize
          \item In-Sample Sharpe: $1.35$
          \item Out-of-Sample Sharpe: \textbf{+1.04}
          \item \emph{Diagnosis:} Cross-asset variance forces attention heads to identify macro regimes common across markets.
        \end{itemize}
      \end{block}
    \end{column}
  \end{columns}
  \vspace{0.2cm}
  \begin{figure}[ht]
    \centering
    \includegraphics[width=0.65\linewidth]{final_audit_results.png}
    \caption{Full-period cumulative net performance and out-of-sample audit breakdown.}
  \end{figure}
\end{frame}

\begin{frame}{HPC Implementation: Overcoming the $O(T^2)$ Memory Wall}
  \begin{alertblock}{The CUDA Out-Of-Memory (OOM) Trap}
    Naively pre-generating sliding windows for 46 assets over 30 years with 252-day lookbacks and 20+ features creates 260,000+ full tensors, requiring $>32\text{ GB}$ of RAM/VRAM during batch initialization and immediately crashing cluster compute nodes.
  \end{alertblock}
  \vspace{0.2cm}
  \begin{columns}[t]
    \begin{column}{0.48\textwidth}
      \textbf{The Zero-Copy Lazy Loader Solution}
      \begin{itemize}\scriptsize
        \item Store flat 1D contiguous arrays in memory.
        \item Utilize NumPy's \texttt{sliding\_window\_view} with zero-copy stride pointers.
        \item Slice sequences on-the-fly inside the PyTorch \texttt{Dataset.\_\_getitem\_\_} method.
        \item Slices memory footprint by \textbf{8x} ($<4\text{ GB}$ VRAM), enabling rapid convergence on standard GPU nodes.
      \end{itemize}
    \end{column}
    \begin{column}{0.48\textwidth}
      \textbf{Cluster Job Submission Practice}
      \begin{itemize}\scriptsize
        \item When submitting to SGE/Slurm clusters (e.g. UCL CS cluster), enforce resource reservations:
        \begin{lstlisting}[language=bash]
# SGE submission flags
#$ -l tmem=16G
#$ -l h_vmem=16G
#$ -l gpu=true
#$ -R y
        \end{lstlisting}
        \item Prevents job starvation and silent kills from unreserved host memory thrashing.
      \end{itemize}
    \end{column}
  \end{columns}
\end{frame}

\begin{frame}{Rigorous Backtesting Protocol: Mitigating Look-Ahead Bias}
  \begin{block}{Three Pillars of Academic Backtest Rigor}
    \begin{enumerate}\scriptsize
      \item \textbf{Point-in-Time Feature Alignment:}
      Signals at bar $t$ must only use information available before close $t$. Return realization occurs at $t+1$:
      \begin{equation}
        \text{PnL}_{t+1} = z_t \cdot r_{t+1} - c \cdot |z_t - z_{t-1}|
      \end{equation}
      Any normalisation using full-sample moments ($\mu_{\text{all}}, \sigma_{\text{all}}$) introduces catastrophic look-ahead bias.
      \item \textbf{Expanding-Window Walk-Forward Retraining:}
      Rather than a static single train/test split, retrain model parameters periodically (e.g., annually):
      \begin{itemize}\tiny
        \item Train on $[1990, 2015] \to$ Test $2016$
        \item Train on $[1990, 2016] \to$ Test $2017$, etc.
      \end{itemize}
      \item \textbf{Multi-Seed Ensembling:}
      Single transformer runs exhibit substantial random seed variance. Wood et al. average the outputs of a 5-seed ensemble to smooth out idiosyncratic execution timing and boost realized Sharpe ratio by $0.4$--$0.8$.
    \end{enumerate}
  \end{block}
\end{frame}

\begin{frame}{Masterclass Synthesis: From Classical Rules to AI Trend Following}
  \begin{columns}[t]
    \begin{column}{0.48\textwidth}
      \textbf{What Survived 200 Years}
      \begin{itemize}\scriptsize
        \item Momentum persistence across macroeconomic regimes.
        \item Positive return skewness and crisis alpha during equity shocks.
        \item The fundamental role of volatility targeting ($1/\sigma$).
      \end{itemize}
      \vspace{0.1cm}
      \textbf{What AI Transformed}
      \begin{itemize}\scriptsize
        \item Replaced rigid moving average parameters with dynamic attention mechanisms.
        \item Formulated end-to-end differentiable loss directly optimizing Sharpe and Sortino ratios.
        \item Enabled dynamic regime recognition without manual parameter recalibration.
      \end{itemize}
    \end{column}
    \begin{column}{0.48\textwidth}
      \textbf{The Rules of Engagement}
      \begin{itemize}\scriptsize
        \item \textbf{Never abandon inductive bias:} Anchor deep models to proven classical priors ($\alpha$-residual modulation).
        \item \textbf{Prioritize data fidelity:} Eliminate futures roll artifacts using ratio adjustment or institutional continuous feeds.
        \item \textbf{Ensure broad cross-sections:} Deep models require multi-asset breadth for regime regularisation.
      \end{itemize}
      \vspace{0.2cm}
      \begin{center}
        \textbf{Day 2 Final Q\&A and Discussion}
      \end{center}
    \end{column}
  \end{columns}
\end{frame}

\begin{frame}{Academic Literature \& Foundational References}
  \begin{block}{\tiny Primary Academic Bibliography (Clickable DOIs)}
    \begin{itemize}\tiny
      \item \href{https://doi.org/10.1111/j.1540-6261.1993.tb04702.x}{Jegadeesh, N., \& Titman, S. (1993)}. Returns to Buying Winners and Selling Losers. \emph{Journal of Finance}, 48(1), 65--91.
      \item \href{https://people.duke.edu/~dah7/RFS2001.pdf}{Fung, W., \& Hsieh, D. A. (2001)}. The Risk in Hedge Fund Strategies: Theory, and Evidence from Trend Followers. \emph{Review of Financial Studies}, 14(2), 313--341.
      \item \href{https://arxiv.org/abs/physics/0508104}{Bouchaud, J.-P., \& Potters, M. (2005)}. Momentum Trading: More Winners than Losers. \emph{arXiv:physics/0508104}.
      \item \href{https://ssrn.com/abstract=2465623}{Bruder, B., \& Gaussel, N. (2011)}. Risk-Return Analysis of Dynamic Investment Strategies. \emph{SSRN Working Paper 2465623}.
      \item \href{https://doi.org/10.1016/j.jfineco.2011.11.003}{Moskowitz, T. J., Ooi, Y. H., \& Pedersen, L. H. (2012)}. Time Series Momentum. \emph{Journal of Financial Economics}, 104(2), 228--250.
      \item \href{https://www.risk.net/derivatives/structured-products/2194247/momentum-trading-skews-me}{Martin, R., \& Zou, D. (2012)}. Momentum Trading 'Skews Me. \emph{Risk Magazine}.
      \item \href{http://www.thierry-roncalli.com/download/Alternative_Risk_Premia.pdf}{Hamdan, R., Pavlowsky, F., Roncalli, T., \& Zheng, B. (2012)}. Understanding the Momentum Risk Premia. \emph{Lyxor Asset Management Working Paper}.
      \item \href{https://doi.org/10.1111/jofi.12021}{Asness, C. S., Moskowitz, T. J., \& Pedersen, L. H. (2013)}. Value and Momentum Everywhere. \emph{Journal of Finance}, 68(3), 929--985.
      \item \href{https://arxiv.org/abs/1404.7797}{Lemp{\'e}ri{\`e}re, G., et al. (2014)}. Two Centuries of Trend Following. \emph{Journal of Investment Strategies}, 3(3), 41--61.
      \item \href{https://ssrn.com/abstract=2695101}{Baz, J., Rattray, M., Harvey, C. R., Van Hemert, O., \& Hofmann, G. (2015)}. Dissecting Investment Strategies: Momentum, Value, Carry. \emph{Man AHL / SSRN Working Paper}.
      \item \href{https://www.nber.org/papers/w20439}{Daniel, K., \& Moskowitz, T. J. (2016)}. Momentum Crashes. \emph{Journal of Financial Economics}, 122(2), 221--247.
      \item \href{https://arxiv.org/abs/1904.04912}{Lim, B., Zohren, S., \& Roberts, S. (2019)}. Enhancing Time Series Momentum Strategies Using Deep Neural Networks. \emph{Journal of Financial Data Science}, 1(2), 19--38.
      \item \href{https://arxiv.org/abs/2112.08534}{Wood, K., Giegerich, S., Roberts, S., \& Zohren, S. (2022)}. Trading with the Momentum Transformer: An Intelligent and Interpretable Architecture. \emph{arXiv:2112.08534}.
    \end{itemize}
  \end{block}
\end{frame}

\end{document}
